% day3_seg1_navproblem/basicnavprob.tex
%\begin{frame}
%Solutions to the Basic Navigation Problem\\
%Solution in Inertial Space
% \begin{itemize}
% \item In an inertial reference frame, it's no problem, if you 
% have a cluster of inertial sensors and you know your initial 
% position:
% \item Just integrate the acceleration once to get velocity, 
%twice to get
% position.
%\[v(T) = v(0) + \int_{0}^{T}\dfrac{d^2r}{dt^2}dt\]
%\[r(T) = r(0) + \int_{0}^{T}v(t)dt\]
% \item And integrate the angular rate once to get attitude.
%\end{itemize}
%\end{frame}

%\begin{frame}
%Solution in Inertial Space
%\begin{itemize}
%  \item Solution easy in inertial space because ...
%  \begin{itemize}
%	\item inertial instruments sense acceleration and rotational % rate with respect to inertial coordinate frame 
%	\item inertial frame attached to sensor platform
%    \item Note: You still need to rotate accelerations from body %to inertial coordinates prior to integration (strapdown)
%  \end{itemize}
%  \item But we don't necessarily want the answer to our 
%questions expressed in terms of inertial coordinates!
%\end{itemize} 
%\end{frame}

\begin{frame}
\frametitle{Dead Reckoning}
\begin{itemize}
  \item Process of determining new position based on extrapolating from previously known position and velocity.
  \item Velocity = Speed and Direction
  \item Unaided inertial navigator is essentially a dead reckoning computer
\end{itemize} 
\end{frame}





